% =====================================================================
%  Subl - Graduation Project Book
%  Chapter: Backend Service & Desktop Agent
%
%  A standalone, compilable application of docs/spec.md.
%  Build:  pdflatex documentation-chapter.tex   (run twice)
%  Drop the preamble into your main book and keep from \section onward.
% =====================================================================
\documentclass[11pt,a4paper]{article}

\usepackage[T1]{fontenc}
\usepackage[utf8]{inputenc}
\usepackage{lmodern}
\usepackage[margin=2.4cm]{geometry}
\usepackage{xcolor}
\usepackage{listings}
\usepackage[most]{tcolorbox}
\tcbuselibrary{listingsutf8}

% ---- palette ----------------------------------------------------------
\definecolor{codebg}{HTML}{F7F8FA}
\definecolor{codeframe}{HTML}{D9DEE4}
\definecolor{kw}{HTML}{0B5FFF}
\definecolor{str}{HTML}{0A7D38}
\definecolor{cmt}{HTML}{7A8794}
\definecolor{num}{HTML}{9AA4AF}
\definecolor{accent}{HTML}{3578FF}   % Subl blue
\definecolor{banner}{HTML}{1F2933}

% ---- base listing style ----------------------------------------------
\lstdefinestyle{subl}{
  backgroundcolor=\color{codebg},
  basicstyle=\ttfamily\footnotesize\color{black!88},
  keywordstyle=\color{kw}\bfseries,
  stringstyle=\color{str},
  commentstyle=\color{cmt}\itshape,
  numberstyle=\tiny\color{num},
  numbers=left, numbersep=10pt,
  frame=single, rulecolor=\color{codeframe},
  framesep=6pt, xleftmargin=16pt, xrightmargin=4pt,
  breaklines=true, breakatwhitespace=true,
  showstringspaces=false, tabsize=2,
  columns=fullflexible, keepspaces=true, upquote=true,
}
\lstdefinestyle{csharp}{style=subl,language=[Sharp]C,
  morekeywords={record,var,async,await,sealed,internal,Result,Guid}}
\lstdefinestyle{python}{style=subl,language=Python,
  morekeywords={Protocol,Optional}}

% ---- filename banner shown above each listing ------------------------
\newcommand{\codefile}[1]{%
  \par\noindent
  \colorbox{banner}{\makebox[\linewidth][l]{%
    \ttfamily\footnotesize\color{white}\hspace{2pt}\#\ #1}}%
  \par\vspace{-2pt}}

% ---- the Explanation box: keeps prose visibly separate from code -----
\newtcolorbox{explain}{
  enhanced, breakable,
  colback=accent!4, colframe=accent,
  boxrule=0pt, leftrule=3pt,
  arc=2pt, left=10pt, right=8pt, top=6pt, bottom=6pt,
  fonttitle=\bfseries\small\color{accent},
  title={Explanation},
  before skip=4pt, after skip=14pt,
}

\begin{document}

% =====================================================================
\section*{Backend Service}
\addcontentsline{toc}{section}{Backend Service}

The backend is a .NET~8 Web API built on Clean Architecture and organised as
\emph{vertical slices}: each feature lives in its own folder containing a
command, a handler, and a validator. The solution is split into four layers ---
\textbf{Domain} (entities and rules, framework-free), \textbf{Application}
(use cases), \textbf{Infrastructure} (database, e-mail, ML and real-time
adapters), and \textbf{Web.Api} (HTTP endpoints). Dependencies point inward
only: outer layers know the domain, never the reverse.

\begin{figure}[h]\centering
  % \includegraphics[width=.8\linewidth]{figures/backend-layers.pdf}
  \fbox{\parbox{.8\linewidth}{\centering\vspace{1.2cm}
  \textit{Figure placeholder: Web.Api $\rightarrow$ Application $\rightarrow$
  Domain, with Infrastructure plugged in at the edges.}\vspace{1.2cm}}}
  \caption{Backend layers and dependency direction (inward only).}
\end{figure}

% ---- Segment 1: Result -----------------------------------------------
\subsection*{Explicit results instead of exceptions}

The backend never throws exceptions for expected failures; every operation
returns an explicit \texttt{Result} that the caller must inspect.

\codefile{api/src/SharedKernel/Result.cs}
\begin{lstlisting}[style=csharp]
public class Result
{
    public bool IsSuccess { get; }
    public bool IsFailure => !IsSuccess;
    public Error Error { get; }

    public static Result Success() => new(true, Error.None);
    public static Result<TValue> Success<TValue>(TValue value) =>
        new(value, true, Error.None);
    public static Result Failure(Error error) => new(false, error);
    public static Result<TValue> Failure<TValue>(Error error) =>
        new(default, false, error);
}
\end{lstlisting}

\begin{explain}
\texttt{Result} makes success and failure first-class values rather than
control flow. A handler returns \texttt{Result.Failure(error)} instead of
throwing, and the API layer maps the error to the right HTTP status in one
place. This keeps the domain free of HTTP concerns and makes every failure
path explicit and testable.
\end{explain}

% ---- Segment 2: a vertical slice -------------------------------------
\subsection*{The shape of a feature: one vertical slice}

Every use case is three small files. We show one slice in full --- submitting a
stress survey --- because the other sixty features in the codebase follow the
same shape. First the \emph{command} (the request) and its \emph{validator}
(structural rules):

\codefile{api/src/Application/Surveys/SubmitSurveyResponse/SubmitSurveyResponseCommand.cs}
\begin{lstlisting}[style=csharp]
public sealed record SurveyAnswerInput(Guid QuestionId, int Value);

public sealed record SubmitSurveyResponseCommand(
    IReadOnlyList<SurveyAnswerInput> Answers) : ICommand<SurveyResultResponse>;
\end{lstlisting}

\codefile{api/src/Application/Surveys/SubmitSurveyResponse/SubmitSurveyResponseCommandValidator.cs}
\begin{lstlisting}[style=csharp]
RuleFor(c => c.Answers).NotEmpty();
RuleForEach(c => c.Answers).ChildRules(answer =>
{
    answer.RuleFor(a => a.QuestionId).NotEmpty();
    answer.RuleFor(a => a.Value)
          .InclusiveBetween(0, SurveyResponse.MaxAnswerValue);
});
\end{lstlisting}

Then the \emph{handler} carries the business rules and persists the result:

\codefile{api/src/Application/Surveys/SubmitSurveyResponse/SubmitSurveyResponseCommandHandler.cs}
\begin{lstlisting}[style=csharp]
// Every answered question must be a known active question.
if (!answeredSet.IsSubsetOf(activeSet))
    return Result.Failure<SurveyResultResponse>(SurveyErrors.UnknownQuestions);

// Every active question must be answered exactly once.
if (answeredSet.Count != activeSet.Count ||
    request.Answers.Count != activeSet.Count)
    return Result.Failure<SurveyResultResponse>(SurveyErrors.MissingAnswers);

var response = SurveyResponse.Create(currentUserService.UserId, answers);
context.SurveyResponses.Add(response);
await context.SaveChangesAsync(cancellationToken);
return Result.Success(result);
\end{lstlisting}

\begin{explain}
The command is a plain immutable record; the validator enforces \emph{shape}
(non-empty, each answer in range $0$--$4$); the handler enforces \emph{meaning}
(all active questions answered exactly once) and delegates scoring to the
domain via \texttt{SurveyResponse.Create}. Because each slice is self-contained,
a feature can be read, tested, and changed without touching the rest of the
system --- this is the single most repeated pattern in the backend.
\end{explain}

% ---- Segment 3: the cross-system contract ----------------------------
\subsection*{The contract with the desktop agent}

This endpoint is the seam between the three systems: the desktop agent posts a
batch of six keyboard features, the handler forwards them to the Python ML
service, and the resulting stress reading is stored.

\codefile{api/src/Web.Api/Endpoints/Stress/SubmitMetrics.cs}
\begin{lstlisting}[style=csharp]
public sealed record Request(
    double MeanDwell, double MedianFlight, double CvFlight,
    double MeanDelFreq, double MeanTotTime, int NKeys,
    int? DeleteCount, DateTimeOffset? CapturedAt);

app.MapPost("stress-sessions/{sessionId:guid}/metrics", async (
    Guid sessionId, Request request,
    ICommandHandler<SubmitMetricsCommand, SubmitMetricsResponse> handler,
    CancellationToken ct) =>
{
    var command = new SubmitMetricsCommand(sessionId, request.MeanDwell,
        request.MedianFlight, request.CvFlight, request.MeanDelFreq,
        request.MeanTotTime, request.NKeys, request.DeleteCount,
        request.CapturedAt?.UtcDateTime);

    Result<SubmitMetricsResponse> result = await handler.Handle(command, ct);
    return result.Match(Results.Ok, CustomResults.Problem);
})
.RequireAuthorization();
\end{lstlisting}

\begin{explain}
The endpoint is deliberately thin: it only maps the wire format to a command
and translates the \texttt{Result} into HTTP (\texttt{200} on success, a
problem response on failure). The field names match the agent's camelCase JSON
exactly, so the contract between the two systems is visible in one place. All
real work --- calling the ML model and persisting the reading --- lives in the
handler, keeping transport and logic separate.
\end{explain}

% ---- Segment 4: cross-cutting concern --------------------------------
\subsection*{Validation as a cross-cutting decorator}

Validators are never called by hand. A decorator wraps every command handler
and runs validation before the handler executes.

\codefile{api/src/Application/Abstractions/Behaviors/ValidationDecorator.cs}
\begin{lstlisting}[style=csharp]
public async Task<Result<TResponse>> Handle(
    TCommand command, CancellationToken cancellationToken)
{
    ValidationFailure[] failures = await ValidateAsync(command, validators);

    if (failures.Length == 0)
        return await innerHandler.Handle(command, cancellationToken);

    return Result.Failure<TResponse>(CreateValidationError(failures));
}
\end{lstlisting}

\begin{explain}
The decorator implements the same \texttt{ICommandHandler} interface it wraps,
so it is transparent to callers. If validation passes it forwards to the inner
handler; otherwise it short-circuits with a \texttt{Result.Failure}. Validation
thus becomes a cross-cutting concern configured once in the dependency-injection
pipeline, instead of boilerplate repeated in every handler.
\end{explain}

% ---- Segment 5: real-world resilience --------------------------------
\subsection*{Healing abandoned sessions}

A desktop agent can crash or lose connectivity mid-session, leaving a session
that never ends and blocks new ones. A background service reclaims them.

\codefile{api/src/Infrastructure/StressDetection/BackgroundServices/AbandonedSessionCleanupService.cs}
\begin{lstlisting}[style=csharp]
private static readonly TimeSpan AbandonAfter = TimeSpan.FromMinutes(15);

private async Task CleanupOnceAsync(CancellationToken ct)
{
    DateTime cutoff = DateTime.UtcNow - AbandonAfter;
    List<StressSession> stale =
        await sessionRepository.GetStaleActiveAsync(cutoff, ct);

    if (stale.Count == 0) return;

    foreach (StressSession session in stale)
        session.MarkAsAbandoned();

    await unitOfWork.SaveChangesAsync(ct);
}
\end{lstlisting}

\begin{explain}
Running every five minutes, the service finds active or paused sessions whose
last activity is older than fifteen minutes and marks them
\texttt{Abandoned} through a domain method, not a raw database write. This
prevents an ``active session lock'' from blocking a user forever if their agent
dies --- a small but essential piece of real-world resilience.
\end{explain}

% =====================================================================
\section*{Desktop Agent}
\addcontentsline{toc}{section}{Desktop Agent}

The desktop agent is a Python application that follows the \emph{same} Clean
Architecture layout as the backend (\texttt{domain}, \texttt{application},
\texttt{infrastructure}, \texttt{presentation}). It captures keystroke
\emph{timing} --- never the keys themselves --- aggregates six statistical
features per batch window, and submits them to the backend. The same blueprint
in a second language keeps the project coherent end to end.

% ---- Segment 1: the boundary -----------------------------------------
\subsection*{Ports: depending on interfaces, not implementations}

The agent's logic depends only on interfaces. Every boundary it needs is
declared as a Python \texttt{Protocol}.

\codefile{desktop-agent/subl\_agent/application/ports.py}
\begin{lstlisting}[style=python]
class ApiGateway(Protocol):
    def login(self, email: str, password: str) -> Tokens: ...
    def register_device(self, profile: DeviceProfile) -> str: ...
    def start_session(self, device_id: str, notes: Optional[str]) -> str: ...
    def submit_metrics(self, session_id: str, features: dict) -> Reading: ...
\end{lstlisting}

\begin{explain}
\texttt{ApiGateway} is the seam between the agent's logic and the network.
Services depend on this protocol, so the real HTTP client can be swapped for a
fake in tests without touching business code. This mirrors the backend's use of
interfaces and keeps the two systems' architecture identical despite being
written in different languages.
\end{explain}

% ---- Segment 2: the novel core ---------------------------------------
\subsection*{Capturing keystroke dynamics (the core)}

This is the heart of the project. As keys are pressed and released, the monitor
records two timings: \emph{flight} (gap between consecutive key-downs) and
\emph{dwell} (how long a key is held). Implausible values are discarded.

\codefile{desktop-agent/subl\_agent/infrastructure/keyboard/pynput\_monitor.py}
\begin{lstlisting}[style=python]
def _on_press(self, key) -> None:
    now = _now_ms()
    with self._lock:
        if self._last_keydown_ms is not None:
            gap = now - self._last_keydown_ms
            if 0 < gap < self._max_reasonable_flight_ms:
                self._flight_ms.append(gap)        # flight time
        self._last_keydown_ms = now
        if kid not in self._press_times:
            self._press_times[kid] = now
            self._total_keystrokes += 1
            if is_del:
                self._delete_count += 1
\end{lstlisting}

Once a batch window closes, the raw timings are reduced to the six features the
ML model expects:

\codefile{desktop-agent/subl\_agent/infrastructure/keyboard/pynput\_monitor.py}
\begin{lstlisting}[style=python]
mean_dwell    = statistics.mean(self._dwell_ms)
median_flight = statistics.median(self._flight_ms)
cv_flight     = statistics.stdev(self._flight_ms) / mean  # coeff. of variation
mean_del_freq = (self._delete_count / self._total_keystrokes) * 100.0
mean_tot_time = _now_ms() - self._first_keydown_ms

features = {
    "meanDwell": round(mean_dwell, 2),     "medianFlight": round(median_flight, 2),
    "cvFlight": round(cv_flight, 4),        "meanDelFreq": round(mean_del_freq, 2),
    "meanTotTime": round(mean_tot_time, 2), "nKeys": self._total_keystrokes,
}
\end{lstlisting}

\begin{explain}
Only \emph{timing and counts} ever leave the keyboard --- never the characters
typed --- which is what makes the system privacy-preserving. All access to the
shared buffers is guarded by a lock because pynput delivers events on its own
thread. The six features (dwell, flight, flight variability, deletion rate,
total time, key count) are exactly the inputs the stress model was trained on,
so this method defines the contract between agent and model.
\end{explain}

% ---- Segment 3: lifecycle + resilience -------------------------------
\subsection*{The session loop and self-healing}

A background thread sends one batch per interval. If the server reports the
session has expired, the agent transparently starts a new one instead of
stopping.

\codefile{desktop-agent/subl\_agent/application/services/session\_service.py}
\begin{lstlisting}[style=python]
def _sender_loop(self) -> None:
    while not self._stop_event.wait(self._batch_interval_seconds):
        self._heartbeat()              # keep the device "online"
        if self._is_paused:
            continue
        self._send_one_batch()

# inside _send_one_batch, when the server returns 404/409:
if ex.status in (404, 409) and self._auto_restart:
    new_id = self._api.start_session(self._device_id, "auto-restart after expiry")
    self._session_id = new_id          # self-heal, keep monitoring
\end{lstlisting}

\begin{explain}
The sender loop is driven by \texttt{wait()} on a stop-event, so the thread
exits promptly when the user stops. A heartbeat is sent every tick --- even
while paused --- so the backend sees the agent as alive and claimable. The
\texttt{404/409} branch is the client-side mirror of the backend's abandoned-
session cleanup: if a session was reclaimed or expired, the agent quietly
acquires a fresh one rather than failing, so monitoring survives transient
server-side changes.
\end{explain}

% =====================================================================
\section*{Summary}
\addcontentsline{toc}{section}{Summary}

The backend and the desktop agent share a single architectural blueprint
expressed in two languages. Typed contracts --- \texttt{Result} and explicit
errors on the server, \texttt{Protocol} ports on the client --- keep every
layer replaceable and independently testable, while the domain in both systems
stays free of framework and transport concerns. The two also cooperate
defensively: the server reclaims abandoned sessions and the agent self-heals
when its session expires, so monitoring is robust to crashes and lost
connectivity. The next chapter describes the machine-learning service that
turns the six keystroke features defined here into a stress prediction.

\end{document}
